Giant Metrewave Radio Telescope (GMRT), one of the world's largest radio
telescopes at metre wavelengths, is located in western India (latitude:
$19^\circ 6'$\,N, longitude: $74^\circ 3'$\,E). GMRT consists of 30 dish
antennas, each with a diameter of 45.0\,m. A single dish of GMRT was held fixed
with its axis pointing at a zenith angle of $39^\circ 42'$ along the northern
meridian such that a radio point source would pass along a diameter of the
beam, when it is transiting the meridian.

What is the duration $T_{\mathrm{transit}}$ for which this source would be
within the FWHM (full width at half maximum) of the beam of a single GMRT dish
observing at 200\,MHz?

\textbf{Hint:} The FWHM size of the beam of a radio dish operating at a given
frequency corresponds to the angular resolution of the dish. Assume uniform
illumination.
